
\section{Build a simple SED system, that predicts the presence and absence of the ten sound categories in 1.2-second audio segments.} 
\label{sec:Build a simple SED system}

Our starting point for a simple SED system is a Temporal Convolutional Network (TCN) architecture. We chose this architecture specifically for several reasons:
\newline
Unlike Recurrent Neural Networks (RNNs), assuming proper dilation settings, they can model long-term dependencies much more efficiently than RNNs, without needing to resort to recursion. They do not rely on sequential processing, which generally makes them much faster to train and easier to parallelize. They are relatively light-weight architectures, especially when compared to transformer based methods. This also makes them less likely to overfit, which we believe might become problem with some of the rarer classes.

Our network takes in the whole F$\times$D array of an audio file, with F representing the number of frames, and D the number of feature (audio embedding) dimensions. The output of the model is then of shape F$\times$10, one probability for each class per frame.

A more detailed description of our architecture: The input is first passed through a Convolutional Feature Extractor block, whose main job is to downsample the D input features into a more compact representation useful for predicting the different classes. In our case, we end up with a F$\times$128 array. For our actual TCN, we then use 4 residual blocks, with dilations of $2^0, \dots, 2^3$ and kernel size of 3. Throughout our model, we also added several BatchNorm and Dropout layers. A more detailed description can be seen in \hyperref[tab:architecture]{Table \ref{tab:architecture}}.

\begin{table}[h]
\centering
\begin{adjustbox}{width=\textwidth}
\begin{tabular}{|c|c|c|c|}
\hline
\textbf{Layer} & \textbf{Input Shape} & \textbf{Output Shape} & \textbf{Details} \\
\hline
\multicolumn{4}{|c|}{\textbf{Convolutional Feature Extractor}} \\
\hline
Input & $[B, F, D]$ & $[B, F, D]$ & Raw frame-level features \\
\hline
Transpose & $[B, D, F]$ & $[B, D, F]$ & Prepare for Conv1D over embedding axis \\
\hline
Conv1D + BN + ReLU & $[B, D, F]$ & $[B, C, F]$ & $1^{\text{st}}$ Conv1D block, kernel size = 3, padding = 1 \\
\hline
Dropout & $[B, C, F]$ & $[B, C, F]$ & Dropout rate = 0.2 \\
\hline
Conv1D + BN + ReLU & $[B, C, F]$ & $[B, C, F]$ & $2^{\text{nd}}$ Conv1D block, kernel size = 3, padding = 1 \\
\hline
\multicolumn{4}{|c|}{\textbf{TCN (Temporal Convolutional Network)}} \\
\hline
ResidualBlock (× num\_blocks) & $[B, C, T]$ & $[B, C, T]$ & Each block: 2 × Conv1D + BN + ReLU + Dropout = 0.2, dilations = $2^i$ \\
\hline
Conv1D (output layer) & $[B, C, T]$ & $[B, N, T]$ & Kernel size = 1 \\
\hline
Transpose & $[B, N, T]$ & $[B, T, N]$ & Final output for framewise classification \\
\hline
\end{tabular}
\end{adjustbox}
\captionsetup{justification=centering}
\caption{Architecture of our network, \\ B: BatchSize, F: Number of Frames, D: Input Dimensionality, C: Reduced Dimensionality, N: Number of Classes}
\label{tab:architecture}
\end{table}

Our model is trained on the training data. After each epoch, it's performance on the validation data is recorded to keep track of and save the best version of the model. We use Binary Cross Entropy as our loss function, a batch size of 64, a learning rate of 1e-4 as well as the Adam optimizer (these setting were kept from the tutorial notebook).

We want to take a moment to point out that our model is pretty lightweight with around 750K parameters, especially when compared to something like the BiGRU used in the tutorial, which has around 20 Million and performs worse. This is of course not part of the task, but noteworthy anyways and should be considered when implementing a system like this.

\subsection{How did you threshold and combine predictions to derive a label for the 1.2-second segments? Describe any heuristics or post-processing strategies used.}
\label{sec:Build:a}
Of course, frame wise probabilities for each class are not enough to classify something. In order to actually predict if a class is present within a given frame we need to binarize our predictions by thresholding them. For now, we decided to use a simple constant threshold value of 0.5 for this.
\newline
Then there is the problem of combining the predictions made for all frames into ones for 1.2 second-segments. Since our hop length is 120ms, 10 frames correspond to 1.2 seconds. Our simple aggregation strategy now is the following: If a class is active within at least one of these 10 frames, the class is said to be active within the segment, otherwise it is labeled as inactive. Since this is in line with how the actual labels are computed, we decided that it most likely would not make a lot of sense to do something different.


\subsection{What strategies did you apply to minimize the task-specific cost function?}
\label{sec:Build:b}
In order to minimize our task specific cost function, which assigns different costs to False Negatives (FN) and False Positives (FP) for each class, we created a custom Binary-Cross-Entropy (BCE) Loss function. It computes the regular BCE Loss for each class, and then by looking at the targets, scales each entry by either it's FP (if the target is 1) or FN (if the target is 0) cost. Doing this improved performance by quite a bit.
\newline
Using this custom loss already lowered the lowest cost on the validation data from 43.17488 to 40.1256.




\subsection{Does your classifier-based SED system achieve lower costs compared to the naive baseline system you established in step 1? If not, describe possible issues.}
\label{sec:Build:c}
Yes, our simple SED system does outperform the naive baseline by quit a bit when it comes to our specific cost function (and likely most other metrics as well). This is not really impressive however, as the costs for False Negatives are higher than for False Positives, so most methods which at least sometimes predict a positive label should perform better.
\newline
In the current state, our SED system achieves a total average cost of about 40.1256 on the validation data.
